\bookchapter{Chapter 19: Troubleshooting and Self-Rescue}

\chapterintro{Phones, like people, occasionally “get sick”. Lagging, crashing, black screen, won't charge... most of these issues aren't that scary, and you can try fixing them yourself before sending it for repair. This chapter is your “phone emergency manual”—try low-cost methods first, then seek a professional doctor if needed.}

\booksection{Phone Lagging and Overheating—The Most Frequently Asked Question}

Honestly, it's totally normal for a phone to start lagging after about a year of use. The reason isn't necessarily that the phone is broken—more often it's “dust accumulation”—junk files, background processes, storage space slowly filling up.

\booksubsection{Diagnose First, Don't Rush to Fix:}

\begin{enumerate}
  \item \textbf{Check storage space}: Go to Settings Storage to see remaining space. If it's under 10\%, the system will become very laggy because it's struggling to free up space. My experience: keep at least 15\%-20\% free space for the phone to run smoothly.
  \item \textbf{Check power consumption ranking}: Go to battery settings to see the power consumption ranking over the past 24 hours. If an app takes a large share but you barely used it, it's running in the background. Uninstall it directly or restrict background activity.
  \item \textbf{Think about what you did recently}: Did it start lagging after a system upgrade? (It's normal to lag for the first few days after a major version update, as the system rebuilds indexes in the background.) Or did it start lagging after installing a certain new app? (Try uninstalling it.)
  \item \textbf{Look at the usage environment}: If the phone is left in the blazing sun, or in a thick case while charging and playing games, overheating is inevitable. High temperature also causes the CPU to automatically throttle, making the phone feel suddenly laggy. Let it cool down and try again.
\end{enumerate}

\textbf{Treatment methods (from simple to complex)}: Restart the phone (solves 60\% of problems) Clean up storage space (delete useless apps and videos) If a specific app is causing trouble, clear its cache or uninstall and reinstall Last resort: Back up data and factory reset.

\booksection{App Crashes and Unresponsiveness—When Apps Misbehave}

An app crashes when you open it, or suddenly freezes while in use. The handling process is like a doctor's consultation—step by step:

\begin{enumerate}
  \item \textbf{Close and reopen first}: Swipe away the app from the background, then reopen it. Most temporary bugs are fixed with this one move.
  \item \textbf{Update the app}: Check the app store for updates. Sometimes it's because the app version is too old and incompatible with the system.
  \item \textbf{Check permissions}: Some apps need a certain permission to work properly, but you didn't grant it. For example, a camera app crashes when opened because it doesn't have camera permission. Go to Settings, find that app, and grant the necessary permissions.
  \item \textbf{Clear cache}: Go to Settings App Management Find that app Storage Clear Cache. Note: "Clear Cache", not "Clear Data"! Clearing data will delete all your login status and settings in the app, equivalent to factory resetting the app.
  \item \textbf{Uninstall and reinstall}: The last resort. Before uninstalling, confirm whether the app's data will be lost—for WeChat, back up chat history first; for games, confirm whether your account is bound to a phone number.
\end{enumerate}

\booksection{Won't Charge or Charges Slowly—It Might Not Be the Battery's Fault}

Not being able to charge is a particularly anxiety-inducing problem, but often the cause is with the accessories, not the phone being broken.

\begin{itemize}
  \item \textbf{Check the charging port first}: Use a toothpick or non-conductive small tool to gently clean out dust and lint from the charging port. I've encountered this myself—a thread from my pocket happened to block the charging port. When I plugged it in, I thought it was poor contact, but after cleaning it out, it was fine.
  \item \textbf{Try a different cable and charger}: Cables have a much higher failure rate than chargers, especially cables that are often rolled up. Swap in a known-good cable and charger for cross-testing to quickly determine which accessory is broken.
  \item \textbf{Wireless charging issues}: If you use wireless charging, check whether the phone case is too thick, or whether there's a metal plate on the back of the case (like a magnetic mount). Also, place the phone precisely in the center of the charging pad—if it's off, it may not charge.
  \item \textbf{Slow charging}: Confirm the charger's power is sufficient (iPhone needs at least a 20W PD charger for fast charging), the cable supports fast charging (not all cables can handle 3A+ current), and the phone isn't too hot (too hot will automatically reduce power to protect the battery).
\end{itemize}

\booksection{Can't Connect to Network, Wi-Fi Keeps Dropping}

Mobile data works fine outside, but can't connect to Wi-Fi at home; or Wi-Fi is connected but there's no internet. Troubleshoot in this order:

\begin{enumerate}
  \item \textbf{Toggle Airplane Mode}: Swipe down to Control Center, turn on Airplane Mode, wait 5 seconds, then turn it off. This is equivalent to making the phone search for the network again, very effective.
  \item \textbf{"Forget" Wi-Fi and reconnect}: Go to Wi-Fi settings, tap the exclamation mark next to the network name (gear icon on Android), select "Ignore This Network" or "Forget", then re-enter the password to connect.
  \item \textbf{Restart the router}: Many people's home routers haven't been turned off for months, and over time they become unstable. Unplug the power, wait 30 seconds, then plug it back in—most Wi-Fi problems are solved.
  \item \textbf{Check if it's the phone's problem}: Do other devices connect to the same Wi-Fi normally? If they all work fine but your phone doesn't, the phone may need a restart or network settings reset (Settings General Transfer or Reset Reset Network Settings—won't delete your data, but will clear all Wi-Fi passwords).
\end{enumerate}

\booksection{Screen Issues: Unresponsive Touch, Flickering, Image Retention}

\begin{itemize}
  \item \textbf{Unresponsive touch}: First peel off the tempered glass film to see if it's the film's fault—cheap tempered films can indeed affect touch over time. Water on the screen or oil on your hands can also cause touch drift. If neither, restart and try again.
  \item \textbf{Marks/ghost image on screen (burn-in)}: Displaying the same image for a long time (like keeping the screen on while browsing TikTok) will leave "ghost image" on OLED screens. This isn't screen damage, just temporary. Turning off the screen and resting for half an hour usually recovers it.
  \item \textbf{Screen flickering}: If it flickers at low brightness, it may be a PWM dimming issue, more common on OLED screens. Look for "DC dimming" or "anti-flicker" option in Settings—turning it on helps a lot. Not all phones have this feature.
  \item \textbf{A completely dark/bright area on the screen}: This is more serious—likely physical screen damage, requiring screen replacement. Non-official repair screens usually have inferior color and brightness compared to originals, and also affect water resistance.
\end{itemize}

\booksection{Won't Turn On, Stuck at Logo—Don't Panic First}

This situation is most likely to make your heart race, but stay calm and try these first:

\begin{itemize}
  \item \textbf{Charge for half an hour first}: Often the battery is just completely drained. Plug in the original charger, wait at least 30 minutes, then try to turn on. Some phones need to charge to a certain level before responding after being fully drained.
  \item \textbf{Force restart}: For iPhone full-screen models, press volume up once, volume down once, then press and hold the side power button until the Apple Logo appears (may need to hold for 10-15 seconds). General Android method: press and hold the power button for 10+ seconds until vibration or screen lights up.
  \item \textbf{Stuck at boot screen, rebooting repeatedly}: This is called "white Apple" (iPhone) or "stuck at Logo" (Android). If force restart doesn't work, you may need to enter recovery mode and use a computer to repair. Connect iPhone to computer, open iTunes/Finder, enter recovery mode, then select "Update" (not "Restore")—this will reinstall the system but preserve data. For Android, enter Fastboot or Recovery mode; key combinations vary by brand, recommend searching your phone model + "enter recovery mode".
  \item \textbf{If none of the above work}: Go to Wi-Fi settings, "Forget" your current Wi-Fi, re-enter password to connect.
\end{itemize}

\booksection{Sound Issues: Speaker Not Working, Low Call Volume}

\begin{itemize}
  \item \textbf{No speaker sound}: First confirm it's not silent mode (iPhone side silent switch), Do Not Disturb mode, or volume turned all the way down. Sounds silly, but this accounts for 40\% of cases. Then check whether Bluetooth is connected to a headphone or speaker.
  \item \textbf{Caller can't hear you}: Check whether the microphone hole is blocked—the small hole at the bottom of the phone is easily blocked by dust. Clean gently with a soft brush. Don't poke hard with a needle or toothpick, you'll damage the dust mesh.
  \item \textbf{Distorted/cracking sound}: If caused by water in the speaker—don't use a hair dryer on hot! Use a paper towel to absorb surface moisture, then let it air dry naturally. Online suggestions of water-ejection apps (playing low-frequency sound to vibrate water out) have some effect, you can try.
\end{itemize}

\booksection{Finding Official Repair, Avoiding Pitfalls}

If you've tried all the above methods and it still doesn't work, you need professional repair. But the repair market is full of pitfalls—let me share my own experiences and lessons learned:

\begin{itemize}
  \item \textbf{Find real official, not "looks official"}: Go to the brand's official website (Apple.com.cn, Huawei.com.cn, etc.) to find the "authorized service provider" list. Non-authorized shops on the street with an Apple logo may not use original parts—after repair, your water resistance is gone in a flash.
  \item \textbf{Must-do before sending for repair}: Back up all data (don't assume data won't be lost during repair), log out of all important accounts (Apple ID, WeChat, Alipay), turn off "Find My Phone". Some repair shops will ask for your lock screen password, which has privacy risks—you'd better log out of accounts yourself before sending it over.
  \item \textbf{Screen replacement pitfalls}: Non-official screens come in "original-pressed" (original glass replaced), "aftermarket-pressed" (non-original glass), and "assembled" (all non-original). Prices differ by several times, but display quality and durability also differ a lot. If your phone is still under warranty, stick with official repair.
  \item \textbf{Battery replacement pitfalls}: Non-original batteries may have falsely labeled capacity, and after half a year of use, degrade more than original batteries do in two years. Apple's official battery replacement price isn't much more expensive than non-official—recommend going official.
  \item \textbf{Data recovery pitfalls}: If your phone gets water-damaged or broken, and the repair shop says "data recovery costs extra", be careful—some unscrupulous merchants will deliberately exaggerate the damage to charge more. Legitimate data recovery quotes a price first, then operates, and doesn't charge if recovery fails.
\end{itemize}

\begin{tcbwarn}{The Universal First Move: Restart Your Phone}
No matter what problem you encounter—lagging, crashing, can't connect to network, sound issues—first restart your phone. This sounds like a brush-off, but in my own experience, at least 60\% of minor issues are resolved after a restart. Phones, like computers, accumulate all kinds of junk in memory over time—restart clears it and starts fresh. So next time your phone has issues, don't rush to search “what to do”—restart first, then see.

\end{tcbwarn}

\begin{tcbpractice}{Memorize Now: These Lifesaving Operations Should Be Second Nature}
\begin{enumerate}
  \item \textbf{Force restart}: Full-screen iPhone: press volume up volume down long press power button until Logo appears. Android: long press power button for 10-15 seconds until vibration.
  \item \textbf{Clean charging port}: Find a small tool (toothpick or SIM ejector pin), gently clean dust from the charging port.
  \item \textbf{Wi-Fi reconnect}: Go to Wi-Fi settings, "Forget" your current Wi-Fi, re-enter password to connect.
  \item \textbf{Toggle Airplane Mode}: Swipe down to Control Center, turn on Airplane Mode, count to 5 in your head, turn it off, see if network recovers.
  \item Screenshot the steps for these operations, or write them in Notes. When problems arise, you don't need to frantically search—just flip through Notes.
\end{enumerate}

\end{tcbpractice}
